\documentclass[10pt,a4paper]{article}
\usepackage[margin=2cm]{geometry}
\usepackage{fontspec}
\setmainfont{Noto Serif}[Scale=0.95]
\setsansfont{Noto Sans}
\setmonofont{Noto Sans Mono}[Scale=0.85]
\usepackage{microtype}
\usepackage{underscore}
\usepackage{graphicx}
\graphicspath{{./}{../../reports/figures/}}
\usepackage{booktabs,tabularx,array}
\usepackage[table]{xcolor}
\usepackage{titlesec}
\usepackage{enumitem}
\usepackage[most]{tcolorbox}
\usepackage[font=small,labelfont=bf]{caption}
\usepackage{float}
\usepackage[hidelinks]{hyperref}
\usepackage{fancyhdr}

\definecolor{accent}{HTML}{1F4E79}
\definecolor{good}{HTML}{2E7D32}
\definecolor{bad}{HTML}{C62828}
\definecolor{soft}{HTML}{F3F6FA}
\titleformat{\section}{\sffamily\Large\bfseries\color{accent}}{\thesection}{0.6em}{}
\titleformat{\subsection}{\sffamily\large\bfseries}{\thesubsection}{0.5em}{}
\titleformat{\paragraph}[runin]{\sffamily\bfseries}{}{}{}[.]
\titlespacing*{\section}{0pt}{1.4em}{0.6em}
\titlespacing*{\subsection}{0pt}{1.1em}{0.4em}
\setlist{itemsep=2pt,topsep=3pt,leftmargin=1.3em}
\setlength{\parindent}{0pt}
\setlength{\parskip}{0.5em}
\renewcommand{\arraystretch}{1.18}
\newcolumntype{Y}{>{\raggedright\arraybackslash}X}
\pagestyle{fancy}\fancyhf{}
\fancyhead[L]{\sffamily\small Synthetic CXR augmentation: progress digest}
\fancyhead[R]{\sffamily\small 9--14 October 2026}
\fancyfoot[C]{\sffamily\small\thepage}
\renewcommand{\headrulewidth}{0.3pt}

\newtcolorbox{keybox}[1][]{colback=soft,colframe=accent,boxrule=0.6pt,arc=2pt,left=8pt,right=8pt,top=6pt,bottom=6pt,
  fonttitle=\sffamily\bfseries,title=#1}
\newtcolorbox{fixbox}{colback=red!3,colframe=bad!70,boxrule=0.5pt,arc=2pt,left=8pt,right=8pt,top=5pt,bottom=5pt}
\newcommand{\fig}[3][\linewidth]{\begin{figure}[H]\centering\includegraphics[width=#1]{#2}\caption{#3}\end{figure}}
\newcommand{\down}[1]{\textcolor{bad}{#1}}
\newcommand{\up}[1]{\textcolor{good}{#1}}

\begin{document}

\begin{center}
{\sffamily\LARGE\bfseries\color{accent} Evaluating synthetic chest X-ray augmentation\\[2pt] for cross-dataset pneumonia classification}\\[6pt]
{\sffamily\large Progress digest: everything that happened, in order, with the numbers that matter}\\[3pt]
{\small Prediction component · state of 14 October 2026 · source of truth: \texttt{docs/FINDINGS_SYNTHETIC.md}, \texttt{docs/EXPERIMENTS.md}}
\end{center}

\begin{keybox}[The short version]
\textbf{Question.} Does generative augmentation (synthetic X-rays from the teammate's Stable Diffusion + LoRA generator) improve a
pneumonia classifier trained on Kermany (children, one hospital) when it is tested on RSNA (adults, another hospital), compared with
no augmentation, conventional augmentation, and the same number of extra \emph{real} images?

\textbf{Answer: no.} In every setting tested, adding synthetic images made the external result \emph{worse}:
\begin{itemize}
\item all 7 architectures, \down{14 of 14} comparisons below zero (\down{$-$0.009 to $-$0.067} AUROC pooled, \down{$-$0.02 to $-$0.15} view-balanced);
\item the same number of \emph{real} images does not hurt (and helps under the fairer metric); simple augmentation gains \up{+0.02 to +0.09};
\item no fix tried (more synthetic images, equalised preprocessing, quality-aware selection, robust losses, augmentation on top)
turns the penalty positive. The best fix makes it small, never a gain.
\end{itemize}
\textbf{Why.} Kermany has an image-quality shortcut (pneumonia films are blurrier and flatter than normal films). The generator
copies that cue and amplifies the contrast part. Every synthetic image is one more lesson in ``blurry = pneumonia'', which is
nearly useless on RSNA and costs most on blurry bedside (AP) films of sick patients.

\textbf{Also found on the way:} (1) RSNA's view (AP vs PA) is strongly confounded with the label; view alone scores 0.796
AUROC, so the usual pooled AUROC flatters every model. (2) A quality-consistency loss (\emph{qcons}) is a useful loss-level
augmentation (+0.031 pooled, +0.066 view-balanced), but it does not make synthetic data help.

\textbf{Needs you:} decide whether the report uses view-balanced AUROC as the headline metric (recommended); write report §4.5;
get three answers from the teammate (§9).
\end{keybox}

\subsection*{What happened when}
\begin{center}\small
\begin{tabularx}{\linewidth}{@{}l Y@{}}
\toprule
\textbf{9--10 Oct} & Data pipeline; Kermany leakage (170 patients) and an acquisition shortcut (file size alone gives 0.998 AUROC);
HPO for 7 architectures; baselines E1 and conventional augmentation E2 (+0.02 to +0.05 for all 7); real-data controls and scaling;
anti-overfitting variants; notebooks 01--06; demo; report §1--§5 drafted. Synthetic images still missing. \\
\textbf{12 Oct} & Teammate's synthetic images arrive. E3 / E4 hurt. Diagnosis: they carry Kermany's image-quality shortcut (§5).
Equalised preprocessing tried; RSNA split into dev / test; synthetic images on top of augmentation (still hurt); dose series
(more synthetic images do not help). \\
\textbf{13 Oct} & Equalisation ``rescue'' withdrawn after re-tuning; shortcut measured in untouched RSNA films; penalty confirmed on
all 7 architectures; quality-aware selection; four robust losses (qcons works, nothing rescues synthetic data); AP/PA view
confound found and quantified; experiment inventory. \\
\textbf{14 Oct} & Gap-filling queue running (missing architecture combinations, §9); notebook 04 refreshed; this digest. \\
\bottomrule
\end{tabularx}
\end{center}

\fig{fig_scoreboard.png}{The whole study in one picture (EfficientNet-B0, RSNA-test). Bars: view-balanced AUROC (mean of the AP and
PA AUROC); ticks: the usual pooled AUROC. Green and blue (real data, augmentation) lie at or right of the dotted baseline; red
(synthetic) lies left of it. Orange are synthetic images plus a fix: the fixes help, but each orange bar is still below the same
recipe without synthetic images (e.g.\ E2_ERASE + qcons + synthetic 0.837 vs E2_ERASE 0.872).}

\newpage
\section{Setup in one page}

\paragraph{Data} \textbf{Kermany} (Guangzhou, children 1--5 y, 5,856 images) is the only training source: a balanced subset of
1,000 real images (500 per class, one per patient) is the default training set. \textbf{RSNA} (NIH, adults, 1024² DICOM) is the
external test: \emph{Lung Opacity} vs \emph{Normal}. RSNA is split by patient into \textbf{RSNA-dev} (30\%, only used once, to
choose the augmentation recipe for the follow-up experiments) and \textbf{RSNA-test} (70\%, every reported number).

\paragraph{Models and protocol} Seven ImageNet-pretrained backbones (ResNet-18/50, DenseNet-121, EfficientNet-B0, ConvNeXt-T,
ViT-S/16, Swin-T). Optuna (25 trials each) tuned on Kermany validation BCE, then frozen. 1,500-step budget, early stopping on
validation BCE, 3 seeds per configuration (368+ runs in total on the RTX 3060 Ti). Effects are reported as the AUROC change of
the 3-seed ensemble with a paired, patient-clustered bootstrap 95\% CI. Differences under about 0.02 are treated as inconclusive
because seed-to-seed variation is that large.

\paragraph{Names used throughout}
\begin{center}\small
\begin{tabularx}{\linewidth}{@{}l Y@{}}
\toprule
E1 & baseline: 1,000 real images, no augmentation \\
R25 / R50 & +250 / +500 \emph{real} images (the matched control for E3 / E4) \\
N250 … N2000, FULL & real-data scaling: 250 to 2,000 real images; all 4,093 \\
E2 & conventional augmentation (rotation, shift, scale, brightness, contrast, gamma) \\
E2_ERASE & E2 + random erasing (best recipe on RSNA-dev for 6 of 7 architectures) \\
E2_ACQ / E2_MIX / E2_ROBUST / LP & acquisition jitter / mixup / everything / frozen backbone (linear probe) \\
E3 / E4 & +250 / +500 \emph{synthetic} images (tag \texttt{sd_lora}); SYN1000, SYN2000, SYN_ONLY for the dose series \\
ERASE_SYN / ERASE_REAL & E2_ERASE augmentation + synthetic / + real images \\
HARM / HARMT & histogram-equalised preprocessing for train and test; HARMT with re-tuned hyper-parameters \\
QM / QC & synthetic images chosen to be quality-matched / quality-counter (§6.4) \\
L\emph{QCONS}, L\emph{QDRO}, L\emph{FOCAL}, L\emph{GCE} & training-loss variants (§6.5) \\
\bottomrule
\end{tabularx}
\end{center}

\paragraph{Three words to keep apart} \emph{Augmentation} = random changes to the same real images during training, no new
images. \emph{Synthetic data} = new images from the generator, added as extra training examples. \emph{Preprocessing} = one fixed
transform of every image, training and test.

\paragraph{Metric} Pooled AUROC on RSNA was the pre-registered endpoint. §7 shows it is inflated by the view; the
\textbf{view-balanced AUROC} (mean of the AUROC within AP films and within PA films) is the fairer number and is shown next to it
where it matters.

\section{Days 1--2: data, and a shortcut in Kermany (9--10 Oct)}

Building patient-level splits turned up two problems in Kermany. \textbf{Leakage:} 170 pneumonia patient IDs appear in both the
official train and test folders; their 457 training images were removed. \textbf{An acquisition shortcut:} normal images are
much larger (median width 1,654 vs 1,120 px) and saved at much higher JPEG quality. A logistic regression on width, height, aspect
ratio and file size, \emph{without a single pixel}, reaches 0.998 AUROC on Kermany validation; width alone 0.95 on Kermany test.

\fig{a_acquisition_by_class.png}{Kermany acquisition properties by class. Pneumonia images are narrower and taller, so the label is
predictable from the file alone.}

Consequences that shaped everything after: images are resized with their full field of view (no letterbox padding that would
encode the size); validation AUROC is useless for model selection (every model reaches about 0.99 after 50 steps), so validation
BCE is used; and the Kermany test set and RSNA measure different things. This shortcut turns out to be the root of the synthetic-data
story: §5 shows its \emph{image-quality} traces survive resizing.

\section{Baselines, augmentation and more real data (10 Oct)}

\paragraph{E1 across 7 architectures} Classic CNNs are the strongest and most stable at 1,000 images (RSNA 0.86--0.88 pooled);
ConvNeXt and the transformers are weaker (0.82--0.84) and unstable. Every model loses 0.03--0.08 from Kermany test to RSNA and
over-calls pneumonia on RSNA (sensitivity 0.84--0.91, specificity 0.56--0.74).

\paragraph{E2 conventional augmentation helps all 7} +0.022 to +0.052 AUROC on RSNA, every CI above zero, with larger gains within
each view than pooled. In-domain validation would have \emph{rejected} it (worse validation BCE for 5 of 7): source validation is
a poor guide to external performance.

\fig{e2_vs_e1_delta.png}{E2 vs E1 per architecture and test set. Augmentation improves RSNA for every architecture.}

\paragraph{More real data} For EfficientNet-B0, augmentation on 1,000 images beats 4$\times$ more real data (E2 +0.056 vs FULL
+0.019) and real data on top of augmentation adds nothing. ViT-S is data-hungry (+0.03 per doubling). This set the bar for the
synthetic images: R25 / R50 (+250 / +500 real) gave about 0 for EfficientNet-B0 and about +0.025 for ViT-S.

\fig{real_data_scaling.png}{Real-data scaling curve. Kermany test keeps moving; RSNA flattens for the CNN.}

\paragraph{Anti-overfitting variants (notebook 06)} Random erasing (E2_ERASE) is the best variant on RSNA-dev for 6 of 7
architectures (Swin-T prefers acquisition jitter); mixup gives the best calibration; a frozen backbone is 0.07--0.15 worse except
for Swin-T. RSNA-dev and RSNA-test agree within 0.01, so this choice is not over-fitted.

\section{The synthetic images arrive, and they hurt (12 Oct)}

The teammate delivered 4,000 normal + 1,000 pneumonia images (Stable Diffusion 1.5 + LoRA, 512², one fixed prompt per class).
They look plausible to the eye.

\fig{synth_montage.png}{Synthetic (top two rows) and real Kermany (bottom two rows) images. A real-vs-synthetic probe on generic
features still separates them perfectly (AUROC 1.000), but there are no near-copies of training images.}

E3 and E4 add 250 / 500 synthetic images to the 1,000 real ones, with nothing else changed.

\fig[0.72\linewidth]{synth_e3e4_views.png}{EfficientNet-B0, RSNA-test AUROC: pooled, PA films, AP films. Grey = baseline, green = + real,
red = + synthetic; hatched = with E2_ERASE augmentation. Synthetic images lower every panel; the drop is worst on AP films, and with
augmentation it is even larger (AP 0.86 $\rightarrow$ 0.64).}

\fig[0.9\linewidth]{fig_synth_archs.png}{The penalty on all seven architectures (Δ vs E1, RSNA-test). Pooled: $-$0.009 to
$-$0.067, all 14 CIs below zero. View-balanced: roughly twice as large. DenseNet-121 is the most resistant, ResNet-18 and
EfficientNet-B0 the most affected.}

\begin{center}\small
\begin{tabularx}{0.92\linewidth}{@{}Y r r r@{}}
\toprule
EfficientNet-B0, RSNA-test, Δ vs E1 & pooled & view-balanced & AP films \\
\midrule
+250 real (R25) / +500 real (R50) & +0.009 / $-$0.000 & \up{+0.027 / +0.017} & +0.056 / +0.045 \\
+250 synthetic (E3) & \down{$-$0.044} & \down{$-$0.075} & \down{$-$0.100} \\
+500 synthetic (E4) & \down{$-$0.065} & \down{$-$0.103} & \down{$-$0.110} \\
synthetic vs the same number of real (R25$\rightarrow$E3 / R50$\rightarrow$E4) & \down{$-$0.053 / $-$0.065} & \down{$-$0.102 / $-$0.120} & \\
E2_ERASE + 250 / + 500 synthetic (vs E2_ERASE) & \down{$-$0.055 / $-$0.060} & \down{$-$0.145 / $-$0.134} & \down{$-$0.25 / $-$0.20} \\
E2_ERASE + 250 / + 500 real (vs E2_ERASE) & $-$0.004 / +0.010 & & \\
\bottomrule
\end{tabularx}
\end{center}

\section{Why: the synthetic images teach Kermany's quality shortcut}

\paragraph{The labels are partly real} Classifiers trained only on real images recognise the synthetic classes (AUROC 0.75--0.91),
but the more robust the classifier, the less pneumonia it sees in synthetic pneumonia (mean p 0.28--0.54 vs 0.96--0.98 for real
pneumonia). Synthetic ``pneumonia'' is mostly borderline.

\paragraph{The quality cue is copied and amplified} In Kermany, low sharpness alone predicts pneumonia at 0.87 AUROC, low contrast
at 0.75. In the synthetic set: 0.86 and \textbf{0.89}. In RSNA, the same cues are nearly useless (0.50--0.63).

\fig{synth_quality_stats.png}{Sharpness (top) and contrast (bottom) by class. Kermany and synthetic classes separate on image
quality alone; RSNA classes do not.}

\paragraph{Causal test} Blurring or flattening the \emph{same healthy} images raises every model's pneumonia score. Baseline and
+real models react as strongly as synthetic ones to blur, so ``blurry = pneumonia'' is learned from the real Kermany data;
the synthetic images add a \emph{contrast} reliance (response to $\times$0.7 contrast doubles, +0.21 $\rightarrow$ +0.41).
Augmentation only protects against mild blur ($\sigma$ = 1).

\fig{synth_degrade.png}{Rise in mean pneumonia probability of healthy images after degradation, per training recipe.}

\paragraph{Visible in untouched RSNA films} Regressing each model's score on the film's own sharpness and its label gives
``how much opacity one SD of blur is worth''. Baseline 0.36--0.39; +500 synthetic \textbf{0.90}: blur moves an E4 model almost as
much as an opacity does. Augmentation halves the ratio, mainly by learning more pathology. The synthetic penalty concentrates on
the blurriest third of AP films, where E3 is at chance (0.54), and these are exactly the bedside films of the sickest patients.

\fig{synth_quality_ratio.png}{Weight of image quality relative to the weight of the opacity label, on real RSNA-test films
(dotted line = baseline). Synthetic training roughly doubles the sharpness weight and, on PA films, triples the contrast weight.}

\section{Trying to rescue the synthetic images (12--13 Oct)}

\subsection{Equalised preprocessing (and the first correction)}
Histogram equalisation of every image removes most of the synthetic contrast cue. The first round looked like a rescue
(penalty $-$0.010 / +0.002). After re-tuning hyper-parameters for the new preprocessing (HARMT), the baseline itself improved and the
penalty came back: \down{$-$0.030 / $-$0.038}. With augmentation, equalisation makes it worse ($-$0.065 / $-$0.109).

\subsection{More synthetic images}
From 250 to 2,000 synthetic images the penalty stays at $-$0.04 to $-$0.065 with no trend: a fixed bias, not a dilution effect.
A model trained on 2,000 synthetic images \emph{only} reaches 0.88 on Kermany test but \textbf{0.50 on RSNA PA films} (chance).

\fig{synth_dose_response.png}{Dose-response: RSNA-test AUROC vs number of synthetic images added (EfficientNet-B0).}

\subsection{Augmentation on top}
On the strong E2_ERASE baseline the penalty is as large or larger for EfficientNet-B0, ViT-S and DenseNet-121, while the same number
of real images changes nothing. Erasing removes local cues, which leaves the global quality cue as an even easier shortcut.

\subsection{Choosing which synthetic images to add}
Subsets with the quality cue removed (quality-matched, QM) or reversed (quality-counter, QC) cut the pooled penalty to $-$0.012 /
$-$0.014 and remove the extra contrast reliance. Under the view-balanced metric about half the penalty remains ($-$0.04). It never
helps.

\subsection{Robust training losses}
Four losses, each with and without the synthetic images: focal (down-weight easy examples), GCE (robust to label noise), group DRO
over label $\times$ sharpness, and \textbf{qcons}: BCE + a penalty on how much the score changes under random blur and contrast
reduction.

\fig{synth_losses.png}{Left: the loss on real data. Middle: the loss with +250 synthetic. Right: the synthetic penalty under each
loss. Only qcons helps clearly; no loss makes the right panel positive.}

\begin{itemize}
\item \textbf{qcons} is a real finding: +0.031 pooled, +0.116 on AP films, +0.066 view-balanced, and it cuts the sharpness weight
to the level of augmentation. It works on ViT-S (+0.040), EfficientNet-B0, Swin-T (AP +0.079) and DenseNet-121 (+0.005), but stays
below conventional augmentation, and on top of augmentation it adds nothing (Δ 0.000). It is a loss-level substitute for augmentation.
\item With synthetic data, qcons shrinks the penalty (EfficientNet-B0 $-$0.044 $\rightarrow$ $-$0.025; Swin-T $-$0.042
$\rightarrow$ $-$0.009, AP 0.532 $\rightarrow$ 0.713); with augmentation + qcons it is $-$0.016. \textbf{It never turns positive.}
\item focal / DRO: small gains (+0.007 / +0.009); GCE under-fits and with synthetic data one seed stays at chance.
\end{itemize}

\section{A second confound: AP vs PA view in RSNA (13 Oct)}

A side-chat hypothesis, checked on every run: in RSNA, \textbf{75\% of AP films (bedside, sicker patients) show opacity, against
17\% of PA films}. AP films are also blurrier and flatter. So the pooled AUROC partly rewards \emph{recognising the view}, and the
Kermany quality shortcut happens to point the same way.

\fig{fig_view_confound.png}{Left: label prevalence by view. Middle: pooled vs view-balanced AUROC for all 121 run $\times$
architecture configurations; pooled is 0.05--0.19 higher, and a ``view detector'' alone would score 0.796 pooled. Right: for every
baseline, an AP opacity is ranked above a PA normal 94--98\% of the time, but a PA opacity above an AP normal only 27--37\%.}

\begin{itemize}
\item The baseline models are to a large extent view detectors: among normal films alone they rank AP above PA with AUROC 0.84--0.92.
\item \textbf{No conclusion flips, most get sharper.} Ranking of all configurations: Spearman 0.985 between the two metrics. Of 48
key comparisons one changes sign, and it strengthens the story: +500 real images \emph{help} (+0.017 balanced).
\item RSNA-dev selection is unaffected (Spearman 0.978).
\item \textbf{Recommendation:} report view-balanced AUROC as the headline external metric, AP / PA separately, pooled as secondary.
\end{itemize}

\section{Claims that were corrected along the way}

\begin{fixbox}
\small
\begin{tabularx}{\linewidth}{@{}p{0.36\linewidth} Y@{}}
\textbf{Earlier claim} & \textbf{Correction} \\
\midrule
``Equalised preprocessing removes the synthetic penalty'' & came from an un-tuned, weak equalised baseline; after re-tuning (HARMT) two thirds of the penalty is back ($-$0.030 / $-$0.038) \\
E4 penalty $-$0.069 & 2 seeds; with the third seed $-$0.065 \\
``More synthetic hurts more'' & only E4 with one weak seed; the dose series shows a plateau, no trend \\
``Quality-aware selection removes most of the penalty'' & true pooled; only about half under the view-balanced metric \\
``Augmentation shields against synthetic images'' & the opposite: the penalty with augmentation is the largest of all ($-$0.145 balanced) \\
Kermany test as evidence for / against synthetic data & +500 real images move it as much ($-$0.060); dropped as uninformative \\
\end{tabularx}
\end{fixbox}

\section{Where things stand (14 Oct)}

\subsection{Done}
\begin{itemize}
\item 371 complete runs on the box; every one has all 3 seeds evaluated. Inventory: \texttt{docs/EXPERIMENTS.md}.
\item Findings log §0--§13: \texttt{docs/FINDINGS_SYNTHETIC.md}. Report §1--§4.4, §4.6--§4.7, §5 drafted: \texttt{docs/report.md}.
\item Notebooks 01--07 executed without errors; notebook 04 re-run today with the synthetic results (it still showed them as
``pending''). Streamlit demo with exported E1 / E2 / E2_ERASE / E3 models.
\end{itemize}

\subsection{Running now: filling gaps in the design}
Nothing had failed, but some combinations were only run for some architectures. Most importantly, the matched real controls
R25 / R50 existed only for EfficientNet-B0 and ViT-S, so for the other five the claim ``synthetic is worse than the same number of
real images'' had no direct test. \texttt{scripts/queue_gaps.sh} (78 runs, about 7 hours, resumable) adds:
\begin{enumerate}
\item R25 / R50 on ResNet-18/50, DenseNet-121, ConvNeXt-T, Swin-T;
\item qcons on E1 for ResNet-18/50 and ConvNeXt-T, and E3 + qcons for five architectures;
\item E2_ERASE + 250 synthetic / + 250 real on ResNet-18/50, ConvNeXt-T, Swin-T.
\end{enumerate}
When it finishes: the comparison scripts (already extended to all seven architectures) are re-run, notebook 07 gets a new
seven-architecture synthetic-vs-real figure, and the findings and inventory are updated.

\subsection{Open: needs a decision or an answer}
\begin{center}\small
\begin{tabularx}{\linewidth}{@{}l Y@{}}
\toprule
\textbf{You} & \textbf{View-balanced AUROC as the headline metric} in the report (recommended; no conclusion changes, effects get clearer). \\
\textbf{Teammate} & (1) Was the generator trained on all Kermany images or only on the classifier's 1,000 (``strict'' vs ``transfer'')? (2) How many LoRA steps? (3) Were images filtered? \\
\textbf{Report} & §4.5 synthetic results (placeholder); §3 note that RSNA-dev chose the ERASE recipe; research-question wording; §5 objectives 3--5; §6 limitations; §7 scripts; §4.7 for all 7 architectures; a robust-loss section. \\
\textbf{Deliverables} & slides; roles / contribution table. \\
\textbf{Optional} & re-tune hyper-parameters for qcons / GCE; view-balanced columns in notebooks 04 / 06. \\
\bottomrule
\end{tabularx}
\end{center}

\section{Where to find things}
\begin{center}\small
\begin{tabularx}{\linewidth}{@{}l Y@{}}
\toprule
\texttt{docs/FINDINGS_SYNTHETIC.md} & every synthetic-data finding with tables and CIs (§0--§13) \\
\texttt{docs/EXPERIMENTS.md} & inventory: ran / planned / abandoned \\
\texttt{docs/report.md} & the technical report draft \\
\texttt{HANDOFF.md} (end) & latest session state, how to resume the GPU queue \\
\texttt{notebooks/07_synthetic.ipynb} & visual walkthrough of everything in §4--§7 of this digest \\
\texttt{reports/results/synth_probe/*.csv} & result tables behind every number \\
\texttt{docs/digest/} & this digest (\texttt{xelatex digest.tex}; \texttt{make_figures.py} regenerates its three summary figures) \\
\bottomrule
\end{tabularx}
\end{center}

\end{document}
